\section{计算机辅助证明与验证}\label{sec:computational-verification}

若干主要结果把人工论证与有限精确计算结合起来。本节准确说明哪些内容由
计算证明，哪些内容由解析论证证明，以及两层证明如何衔接。

\subsection{证明分类}

Floquet 直和、行列式恒等式、精确正系数展开、算子等价、消去基线、矩阈值判据、从
分组矩恒等式到缺陷不等式的转换，以及手征约化，都是第~
\ref{sec:preliminaries} 节和第~\ref{sec:periodic-construction}--
\ref{sec:eight-barrier} 节给出的代数证明。

下列断言是有限计算机辅助定理：

\begin{enumerate}
\item 每个偶数阶 \(8\leq n\leq30\) 的每个切换类都满足猜想下界；
\item 表示周期至 16 的 2,626 个合法通量/二面体轨道代表元构成完整有限相空间；
\item 定理~\ref{thm:low-period-frontier} 使用的精确证书划分覆盖上述每个代表元；
\item 式~\eqref{eq:general-flux-moments}--\eqref{eq:general-excesses} 背后的系数归集以及所述 \(F_4,F_6,F_9\) 数值，是由一般周期矩检查器复现的精确计算机辅助符号恒等式；
\item 闭游走矩的首个正指标和五个剩余竞争相位的 Rayleigh 商，是整数或有理运算的精确输出。
\end{enumerate}

独立重新执行与独立编写的检查器用于增强可复核性，但不是定理的附加假设。

\subsection{有限判定的精确性}

浮点特征值只用于定位可能的最大化纤维或提出小有理向量。有限判定仅以以下
一种精确形式进入证明：

\begin{itemize}
\item 由整数矩阵和整数向量计算的有理 Rayleigh 不等式；
\item 由无分数消元或有理消元得到的正定性证书；
\item 整系数多项式的 Sturm 根隔离证书；
\item 精确整数闭游走矩超额量；
\item 精确轨道计数与轨道大小之和。
\end{itemize}

所以任何严格定理不等式都不依赖浮点容差。

\subsection{最小性计算}

对 \(n=8,\ldots,20\)，直接枚举全部 \(2^{n+1}\) 个切换类。对
\(n=22,\ldots,30\)，计算枚举规范四边形通量手镯，保留两种和乐，并附上
精确二面体轨道重数。所代表的切换类总数与 \(2^{n+1}\) 核对。因而在最大
阶数上，是用 17,929,600 个规范谱状态代表 2,147,483,648 个切换类，而不是
运行 21.47 亿个独立测试用例。

对每个大阶数 \(n=24,26,28,30\)，规范 \(Q\) 代表元集合还与一份独立编写的
C 检查器逐记录比较。正式路线采用固定重量 FKM 项链递归，其算法背景见
\cite{FredricksenMaiorana1978,Sawada2001}。独立路线则扫描全部 \(2^n\) 个
二进制字，筛选奇偶性，直接构造二面体像，并在位图中标记完整轨道。两条路线
不共享遍历或规范化代码。逐项消耗集合会核对每个规范字及其轨道大小；集合相等
并非由共同摘要或相同基数推断。

对各正式计算阶数执行了新的确定性重新计算：

\begin{table}[tbp]
\centering
\caption{大阶数枚举中的规范谱状态和检查点分块。}
\label{tab:checkpoint-counts}
\begingroup
\renewcommand{\arraystretch}{1.12}
\begin{tabular}{lll}
\toprule
\(n\) & 规范谱状态 & 分块数 \\
\midrule
24 & 353,812 & 28 \\
26 & 1,299,064 & 76 \\
28 & 4,810,472 & 250 \\
30 & 17,929,600 & 908 \\
\bottomrule
\end{tabular}
\endgroup
\end{table}

重新计算得到的有序输入摘要、有序证书摘要、壳层计数、代表总数和末端检查点
序列均与原计算一致，未发现数学判定不一致。已存计算检查点的完整性复核与此次
重新计算的作用不同：前者认证既有执行记录，但不重建每个状态的不等式。

四个阶数的独立代表元数依次为 176,906、649,532、2,405,236 和 8,964,800。
每阶的缺失、重复、未消耗及轨道大小不一致数均为零。独立重构的轨道大小之和
为 \(2^{n-1}\)，添入两种和乐和两个全局取反提升后得到全部 \(2^{n+1}\) 个
切换类。这是覆盖证书，与每个正式谱状态附带的精确判定证书相互分离。

第二条完整判定路线进一步弥合了这种分离。C 扫描器按整数递增顺序导出独立
重构的规范记录；另一份 Python 检查器不导入正式搜索中的赋号、矩阵、阈值或
证书程序，自行从每条记录重构两种 Hamilton 和乐及符号邻接矩阵。浮点特征
向量只用于提出整数向量，随后以整数算术计算
\(\lVert Av\rVert^2/\lVert v\rVert^2\)，并与 \(\rho_-(n)^2\) 的认证有理上端点
比较。特殊最优状态则由特征多项式的精确整除性核验。对
\(n=24,26,28,30\)，该路线依次判定 353,812、1,299,064、4,810,472 和
17,929,600 个谱状态；每阶恰有一个精确最优状态，且没有未认证的非最优状态。

\subsection{低周期计算}

对每个 \(1\leq p\leq16\)，一条路线把所有合法通量字显式划分为二面体轨道，
另一条路线使用 Burnside 引理和置换圈奇偶性。把代表元集合而不只是其基数与
独立生成的规范集合比较；确定性轨道标识符绑定于字典序，并逐行检查合法
奇偶性、规范性、轨道大小、\(\tau\) 提升、本原周期、算子等价和精确证书。

所得 2,626 行表没有缺失或重复的规范代表元。精确排除划分式~
\eqref{eq:frontier-partition} 没有重叠，也不存在未判定行。

\subsection{可复现环境与复核边界}

用于复核的证明材料固定于公开不可变快照

\begin{center}
\url{https://github.com/whzy3185/math/tree/c81be34a3b12a7ac47adbb4499c475df7bf4fc04}
\end{center}

机器可读的投稿清单与强化计算证据清单
\path{research/reproducibility/}\allowbreak\path{target_a_submission_}\allowbreak\path{artifact_manifest.json}
以及
\path{research/reproducibility/}\allowbreak\path{target_a_computational_}\allowbreak\path{evidence_manifest.json}
将定理~\ref{thm:smallest-counterexample} 和定理~
\ref{thm:low-period-frontier} 与检查器源码、固定依赖、24 至 30 阶计算检查点
序列、重新计算摘要和完整 2,626 行低周期表的 SHA-256 摘要对应起来。对定理~
\ref{thm:low-period-frontier}，该表包含每个规范代表元及其精确判定。对定理~
\ref{thm:smallest-counterexample}，已存计算检查点为 24 至 30 阶提供逐块
完整性复核，而确定性搜索程序为每个阶数提供逐状态的完整数学重新计算路线。
直至 22 阶的紧凑材料是汇总证明，不替代该重新计算。

参考环境为 Python 3.12.13、NumPy 2.3.5、SymPy 1.14.0 和 pytest 9.1.1。
机器可读需求文件固定了 Python 包。因而可在干净的 Python 3.12 环境中，
使用附录~\ref{app:computational-protocol} 所列仓库相对命令运行默认套件和三个
显式启用的慢速生成器核验。

仍有两个复核边界；对此类大规模穷举证明，明确披露复核边界是必要做法
\cite{Lam1991}。第一，两条枚举路线虽然使用不同算法、语言和数据结构，却
必然实现同一个已证明的数学商 \((Q,\alpha)/D_n\)。第二，仓库没有把每个大阶
整数 Rayleigh 向量分别保存为单独文件；正式与独立判定路线会重新生成向量并
验证不等式，而各自的有序摘要与检查点序列绑定运行结果。对直至 22 阶的情形，
紧凑历史材料同样认证汇总记录，而非每个状态各存一份可复核文件。这些属于
执行层的复核边界，并非浮点容差或未经核对的轨道计数。
